\scanpage{13}

% ll. 1–3
\mspage{\textcircled{\scriptsize 10 bis}}
d) Le stabilisateur d'un élément de $\mathrm{Pl}_6$ \deleted{dans} \inserted{(ou, ce qui revient} au $=$, de $\mathrm{Pl}_5$ ou $\mathrm{Pl}_7$) dans $W$ est isomorphe à $\unsure{\mathfrak S|\mathcal S}_3$ (via l'opération sur l'ens $I$ \inserted{de cardinal 3} des éléments de

% ll. 4–9
\deleted{$E(u)^*=E(u)-\{v,w\}$} $E(u)^*=E(u)-\{v,w\}$ qui sont $\ne u'$, \deleted{non liés à $\unsure{u'|v'|v}$} et liés à $v'$, i.e. des él. de $E(u)-\{v\deleted{\unsure{,w}},u',\unsure{w}\}$ qui sont liés à $v'$, i.e. $I=E(u)\cap E(v')-\{u',v\}$. On voit \inserted{donc} que cet $I$ est can. is. aux $(9.4):2=18$ ensembles analogues tels que $E(v'')\cap E(w')-\{u'',w'\}$ etc $\dots$ $\}$. Donc $W$ est de cardinal égal à $\operatorname{card}\mathrm{Pl}_6\,.6$ \unsure{$\Rightarrow$} i.e.

% l. 10
\[
\operatorname{Card}W=2^7.3^4.5=51840
\]

% centred rule
\begin{center}\rule{1.15cm}{.35pt}\\[-1.2ex]\rule{.95cm}{.35pt}\end{center}

% below rule, left
\deleted{\unsure{Graphe|Compte}}\\
\deleted{Complexe cubique}
% archive sheet number: 12
